%% ============================================================================
%% Section: Background and Problem
%% Included from main.tex via \input{sections/background}
%% ============================================================================
\section{Background and Problem}\label{sec:bg}

\subsection{Causal DAGs and their property-graph encoding}\label{sec:bg:dag}

A \emph{causal DAG} is a directed acyclic graph $G=(V,E)$ whose vertices $V$ are
random variables and whose edges $E$ encode direct causal influence: an edge
$X\!\rightarrow\!Y$ asserts that $X$ is a direct cause of $Y$ relative to the
other modeled variables~\cite{pearl2009causality}. Variables are either
\emph{observed} or \emph{latent} (unmeasured); latent variables are retained
because they change which conditional independencies hold and which causal
effects are identifiable, even though they are never assigned data. We write
$\mathrm{pa}(X)$, $\mathrm{ch}(X)$, $\mathrm{an}(X)$, and $\mathrm{de}(X)$ for
the parents, children, ancestors, and descendants of $X$; ancestors and
descendants are the transitive closures of parents and children and exclude $X$
itself.

Causal4j stores the DAG as a native property graph. Each variable is a node
\[
  (:\textsf{Variable}\;\{\textsf{id},\ \textsf{name},\ \textsf{observed},\ \textsf{latent}\})
\]
and each direct cause is a relationship
\[
  (:\textsf{Variable})\text{-}[:\textsf{CAUSES}\;\{\textsf{active},\ \textsf{weight}\}]\!\rightarrow\!(:\textsf{Variable}).
\]
The \textsf{id} property is the canonical key; \textsf{active} lets a curator
toggle an edge without deleting it, and only edges with $\textsf{active}=\text{true}$
participate in a model; \textsf{weight} carries an optional structural
coefficient used by the quantitative layer (\S\ref{sec:quant}) and ignored by the
structural queries. This encoding keeps the causal model \emph{in} the database,
co-located and jointly queryable with any other graph data, rather than in a
transient in-memory object owned by an external tool.

\subsection{The causal query catalog}\label{sec:bg:catalog}

Analysts interrogate the DAG through a small, well-established
catalog of structural queries. We group them by the reasoning they support.

\paragraph{Reachability queries} return $\mathrm{pa}$, $\mathrm{ch}$,
$\mathrm{an}$, and $\mathrm{de}$ of a variable, and test whether a directed path
$X\rightarrow\!\cdots\!\rightarrow Y$ exists. They underpin every other query.

\paragraph{Path-role queries} classify intermediate variables on the paths
between a treatment $T$ and an outcome $Y$. A \emph{mediator} lies on a directed
path $T\rightarrow\!\cdots\!\rightarrow Y$; a \emph{confounder} is a common cause,
i.e.\ a variable in $\mathrm{an}(T)\cap\mathrm{an}(Y)$; a \emph{collider} on a
path is a variable with two incoming path edges ($\rightarrow Z \leftarrow$).

\paragraph{Independence queries} test \emph{d-separation}. A path between $X$ and
$Y$ is \emph{blocked} by a conditioning set $Z$ if it contains a non-collider in
$Z$ or a collider none of whose descendants is in $Z$; $X$ and $Y$ are
d-separated given $Z$, written $X \!\perp\! Y \mid Z$, when \emph{every} path
between them is blocked~\cite{pearl2009causality}. D-separation is the graphical
criterion for conditional independence and is \emph{not} plain reachability: a
variable added to $Z$ can block one path while unblocking another through a
collider, so the test must be conditioning-aware.

\paragraph{Identification queries} decide whether the causal effect of $T$ on $Y$
can be estimated from observational data and, if so, return an \emph{adjustment
set}. A \emph{backdoor path} from $T$ to $Y$ is a path with an edge into $T$; a
set $Z$ satisfies the \emph{backdoor criterion} if it blocks every backdoor path
and contains no descendant of $T$. Enumerating minimal valid adjustment sets
issues many d-separation tests, so identification is the most traversal-intensive
query in the catalog.

Table~\ref{tab:operators} (\S\ref{sec:lang}) maps each of these to a Causal4j
language operator and the physical operator that answers it.

\subsection{Running example}\label{sec:bg:example}

We use throughout the COVID-19 causal DAG of Piccininni et
al.~\cite{piccininni2020covid}, comprising $25$ variables and $76$ directed
edges; Fig.~\ref{fig:covid} shows the fragment relevant to one query. Consider
estimating the effect of \textsf{interventions} ($T$) on
\textsf{reported\_cases} ($Y$). The causal (front-door) influence flows along
directed paths such as
$\textsf{interventions}\!\rightarrow\!\textsf{exposure\_sarscov2}\!\rightarrow\!
\textsf{onset\_covid19}\!\rightarrow\!\textsf{reported\_cases}$, whose interior
variables are mediators, and here happen to be \emph{latent}, which the engine
must nonetheless reason about. \textsf{covid19\_burden} is a common cause of both
$T$ and $Y$ (it causes \textsf{interventions} and also
\textsf{exposure\_sarscov2}), so it is a confounder that opens the backdoor path
$\textsf{interventions}\!\leftarrow\!\textsf{covid19\_burden}\!\rightarrow\!
\textsf{exposure\_sarscov2}\!\rightarrow\!\cdots\!\rightarrow\!\textsf{reported\_cases}$.
Conditioning on $\{\textsf{covid19\_burden}\}$ blocks this path and satisfies the
backdoor criterion, so it is a valid adjustment set for $T\!\rightarrow\!Y$. Every
claim in this paragraph is the output of a query in the catalog above.

\begin{figure*}
    \centering
    %% ============================================================================
%% Figure: running example -- fragment of the COVID-19 causal DAG
%% (Piccininni et al. 2020). Edges are faithful to dev/seed-dag.cypher.
%% Included from sections/background.tex inside a figure env.
%% Requires: \usepackage{tikz}; \usetikzlibrary{arrows.meta,positioning}
%% ============================================================================
\begin{tikzpicture}[
    font=\scriptsize,
    >={Stealth[length=1.8mm]},
    node distance=6mm,
    obs/.style={draw, rounded corners=1.5pt, minimum height=5.5mm,
                inner sep=2.5pt, align=center, fill=white},
    lat/.style={draw, dashed, rounded corners=1.5pt, minimum height=5.5mm,
                inner sep=2.5pt, align=center, fill=white},
    treat/.style={obs, fill=black!12, very thick},
    out/.style={obs, fill=black!12, very thick},
    conf/.style={obs, draw=black, densely dotted, very thick},
    edge/.style={->, gray!70},
    bd/.style={->, thick},               % backdoor-path edge
    every node/.style={outer sep=0pt}
]

% --- nodes (x,y placed by topological flow, left -> right) ---
\node[conf]  (burden) at (0.0, 1.5) {covid19\_burden \\ \emph{(confounder)}};
\node[treat] (interv) at (0.0, 0.0) {interventions \\ ($T$)};
\node[obs]   (mob)    at (2.6, 0.8) {mobility};
\node[lat]   (exp)    at (5.0, 0.4) {exposure\_sarscov2};
\node[obs]   (acc)    at (6.4, 1.7) {access\_tests};
\node[lat]   (onset)  at (7.6, 0.4) {onset\_covid19};
\node[out]   (rep)    at (10.2, 0.4) {reported\_cases \\ ($Y$)};

% --- causal (front-door) edges ---
\draw[edge] (interv) -- (mob);
\draw[edge] (burden) -- (mob);
\draw[edge] (mob)    -- (exp);
\draw[edge] (exp)    -- (acc);
\draw[edge] (acc)    -- (onset);

% --- backdoor path: T <- burden -> exposure -> onset -> reported ---
\draw[bd] (burden) -- (interv);
\draw[bd] (burden) to[bend left=12] (exp);
\draw[edge] (interv) to[bend right=12] (exp);
\draw[bd] (exp)    -- (onset);
\draw[bd] (onset)  -- (rep);

\end{tikzpicture}

    \caption{Caption}
    \label{fig:example}
\end{figure*}


\subsection{Problem statement}\label{sec:bg:problem}

We target the following data-engineering problem: \emph{support the causal query
catalog as a first-class workload inside a graph DBMS}. Concretely, a solution
must meet four requirements.
\begin{itemize}
  \item \textbf{R1 (Expressiveness)}. A declarative interface in which each catalog query
  is a single statement, close to how analysts phrase causal questions, that
  compiles to physical operators over the graph store (\S\ref{sec:lang}).
  \item \textbf{R2 (Efficiency)}. Structural queries answered without re-traversing the
  stored graph on every call, so that latency is sublinear in graph size and the
  system sustains the many independence tests an identification query issues
  (\S\ref{sec:index}).
  \item \textbf{R3 (Correctness)}. Results faithful to the graphical semantics above,
  in particular a conditioning-aware d-separation test and a sound backdoor
  criterion (\S\ref{sec:index}, validated in \S\ref{sec:eval}).
  \item \textbf{R4 (Maintainability)}. The accelerating index kept consistent under edge
  and variable updates without full recomputation, since curated causal graphs
  change continually (\S\ref{sec:incr}).
\end{itemize}
None of these is served by exporting the DAG to an external library: R2- R4 in
particular are database concerns, indexing, query processing, and incremental
maintenance, which is why we address them inside the DBMS.
